\section{构造原理}

十九个下界来自扩大受约束构型的不同途径。紧接触并不总会阻止运动：25维的552点层
在移动中保持内部 Gram 矩阵不变。一个不能直接补点的选择，也未必用尽候选几何的空间：
27维的重着色改变了第二层选择，共享阻挡则让一组支持的联合插入产生正收益。
局部容量还可以借助另一种组合对象达到，权四设计到权八符号码的 Hadamard 转移
给出了具体实例。

另外两类机制利用母格提供的信息。格最小范数自动保证38维的跨壳相容性；
设计矩确定八个固定坐标截面的点数，其中包含43维的计数，并将45维的点数归结为
锚图统计量。后一种情形中，完整的298点纤维精确确定7380090点构型。
49至55维则由同一母类产生七个提升，精确并集计数将母类扩充与像族选择的增益分别确定。

Qiushi Engine 通过持续自主数学研究发展这些构造，联系了几何自由度、离散选择和精确计数。
相应证明与有限对象使这些联系成为可以继续研究的数学材料：读者能够重建、改变它们，
并以此发展新的构造。
